\documentclass[10pt,a4paper]{amsart}
\usepackage[margin=19mm]{geometry}
\usepackage{amsmath,amssymb,mathtools}
\usepackage[scr=rsfs,cal=euler]{mathalfa}
\usepackage{xfrac}
\usepackage[unicode,hidelinks,bookmarks=false]{hyperref}
\newcommand{\ie}{i.e.\ }
\newcommand{\cf}{cf.\ }
\newcommand{\resp}{respectively\ }
\newcommand{\Subsection}{\subsection}
\providecommand{\nobreakdash}{\nobreak\hbox{-}}
\setlength{\emergencystretch}{2em}
\multlinegap0pt
\usepackage{bm}
\usepackage{bbm}
\usepackage{dsfont}
\usepackage{xspace}
\usepackage{cancel}
\usepackage{mathtools}
\usepackage{cleveref}
\usepackage[shortlabels]{enumitem}
\usepackage{amsthm}
\makeatletter
\@ifundefined{theorem}{\newtheorem{theorem}{Theorem}}{}
\@ifundefined{corollary}{\newtheorem{corollary}[theorem]{Corollary}}{}
\@ifundefined{lemma}{\newtheorem{lemma}[theorem]{Lemma}}{}
\makeatother
\makeatletter
\newcommand{\NN}{\mathbb{N}}
\newcommand{\PP}{\mathbb{P}}
\newcommand{\Proj}{\operatorname{Proj}}
\newcommand{\cO}{\mathcal{O}}
\newcommand{\fm}{\mathfrak{m}}
\expandafter\ifx\csname lcverbatim\endcsname\relax
\newenvironment{lcverbatim}
 {\SaveVerbatim{cverb}}
\fi
\makeatother
\title[Bertini's theorem for $F$-rationality is false]{Bertini's theorem for $F$-rationality is false}
\author{Thomas Polstra, Austyn Simpson}
\date{Source: arXiv:2609.18921v1 (CC BY 4.0)}
\begin{document}
\maketitle

\noindent\emph{Source and licence.} Bertini's theorem for $F$-rationality is false. Version arXiv:2609.18921v1 (CC BY 4.0), distributed under the Creative Commons Attribution 4.0 International licence, as recorded in the arXiv licence metadata for this exact version and shown on the abstract page https://arxiv.org/abs/2609.18921v1. Full rights notice: licenses/arxiv-2609.18921-CC-BY-4.0.txt.\par\smallskip

\noindent\textit{Editorial scope.} Counted unit: the Lemma whose source label is \texttt{lem: the embedding of the projective variety} in the article (source0.tex lines 945--952, source label \texttt{lem: the embedding of the projective variety}), delivered together with the article's own complete original proof of it (source0.tex lines 954--1032, one \texttt{proof} environment of 79 lines). No mathematics is written, repaired or re-derived: statement and proof are the article's own text line for line. Required context retained uncounted: lem: An affine cover of X (lines 855--862); cor: X is F-rational (lines 868--875). Every definition, display and label of those lines is retained unchanged. Omitted from this package and not redistributed: 1--854, 863--867, 876--944, 953--953, 1033--1145. Nothing inside a retained excerpt is omitted or altered: every retained body is delivered exactly as the article's lines are written, so no omission receipt of any kind accompanies this package. Citations: the retained excerpts carry no citation command at all, so no bibliography entry of the article is transcribed and no reference is redistributed. Reference adaptation: none was needed and none was made; every reference inside the delivered unit resolves inside it, so no unresolved reference or citation is produced. Printed slips kept literal: no printed word, symbol, hypothesis or number of the counted statement or of its proof was altered. Layout and class: the article's own class is \texttt{amsart} with packages including geometry, inputenc, mathptmx, mathalfa, mathrsfs, setspace, amssymb, amsopn, amsmath, array, pdfsync, url, amsfonts, float; the wrapper is the lane's pinned \texttt{amsart} editorial wrapper at 10pt on A4, which restates the theorem-like environments the retained text needs and adds the article's own macro definitions that the retained text needs (\texttt{\textbackslash NN}, \texttt{\textbackslash PP}, \texttt{\textbackslash Proj}, \texttt{\textbackslash cO}, \texttt{\textbackslash fm}), with extra wrapper packages (bm, bbm, dsfont, xspace, cancel, mathtools, cleveref, enumitem). The delivered numbering is the unit's own. Assets: no retained span contains a figure environment, a graphics-inclusion command, a PDF-page inclusion, a table or a TikZ picture; no image file is copied into this package, the package declares no asset dependency and no image file is redistributed with it. Licence: version arXiv:2609.18921v1 of this article is distributed under the Creative Commons Attribution 4.0 International licence, as recorded in the arXiv licence metadata for this exact version and shown on the abstract page https://arxiv.org/abs/2609.18921v1. A full rights notice is delivered at licenses/arxiv-2609.18921-CC-BY-4.0.txt. Characters: every retained excerpt is pure ASCII, so the delivered engine, XeLaTeX with its default Latin Modern text font, reports no missing character.
\begin{lemma}
    \label{lem: An affine cover of X}
    Let $\mathfrak{X} = \Proj_{k}(\mathcal{R})$ and
    \[
    V'=\{T_az^2\mid a\in\Lambda\}\cup \{T_a^3y^3\mid a\in\Lambda\}\cup\{T_a^3x_i^6\mid a\in\Lambda\mbox{ and }1\leq i\leq 3 \}.
    \]
    Then $V'$ is a collection of monomials of $\left(\mathcal{R}\right)_1$ whose associated linear system is base point free.
\end{lemma}

\begin{corollary}
    \label{cor: X is F-rational}
    Let $\mathfrak{X} = \Proj_k(\mathcal{R})$.
    \begin{enumerate}[label=(\alph*)]
        \item For each closed point $x\in \mathfrak{X}$, the embedding dimension $\dim_k(\fm_x/\fm_{x}^2)\leq 20$.
        \item The projective variety $\mathfrak{X}$ is $F$-rational.
    \end{enumerate}
\end{corollary}

\begin{lemma}
    \label{lem: the embedding of the projective variety}
    For each $n\in\NN$, let $\mathcal{E}_n = \{x_1^ax_2^bx_3^cy^d\mid a+b+c+2d = n\}$. Let $\mathfrak{X} = \Proj_k(\mathcal{R})$ and consider the sets
\begin{align*}
    V_1 = \{T_az^2\mid a\in\Lambda\},\quad V_2 = \{T_a^2zm\mid a\in\Lambda\mbox{ and }m \in\mathcal{E}_3\}, \quad V_3 = \{T_a^2T_bm\mid a,b\in\Lambda\mbox{ and }m \in\mathcal{E}_6\}.
\end{align*}
    Then $V := V_1\cup V_2\cup V_3$ is a collection of $2548$ monomials in $\left(\mathcal{R}\right)_1$ that define a very ample linear system.
\end{lemma}

\begin{proof}
    It is clear that each element of $V$ is a monomial of $\left(\mathcal{R}\right)_1$, that $|V_1|=7$, $|V_2| = 7\times13=91$, $|V_3| = 49\times50 = 2450$, and $|V| = 7+91+2450 = 2548$. Let $V'$ be the collection of monomials described in \cref{lem: An affine cover of X}. Then $V'\subseteq V$ and hence $V$ defines a base point free linear system. 

    Let $\varphi_V:\mathfrak{X}\to \PP^{|V|-1}_k = \PP^{2547}_k$ be the associated morphism of $V$. It remains to check $\varphi_V$ is a closed immersion; to do so, it suffices to show that for each $W\in V'$, $\mathcal{O}_\mathfrak{X}(D_+(W)) = k\left[\frac{U}{W}\mid U\in V\right]$. Indeed, for each $W\in V'$ let $Z_W$ be the corresponding coordinate of $\PP^{2547}_k$. Assuming that for each $W\in V'$ that $\mathcal{O}_\mathfrak{X}(D_+(W))$ is generated by the sections $\frac{U}{W}$ as $U$ varies in $V$ implies $D_+(W)\subseteq D_+(Z_W)$ is a closed immersion. Hence $\varphi_V$ defines a closed immersion in the open set $\bigcup_{W\in V'}D_+(Z_W)\subseteq \PP^{2547}_k$. But $\varphi_V:\mathfrak{X}\to \PP^{2547}_k$ is proper, hence its image is closed, so $\varphi_V:\mathfrak{X}\hookrightarrow \PP^{2547}_k$ is a closed immersion as needed.

    It remains to check for every $W\in V'$ the sections $\{\frac{U}{W}\mid U\in V\}$ generate $\cO_{\mathfrak{X}}(D_+(W))$ as a $k$-algebra. First suppose $W = T_az^2$ for some $a\in\Lambda$. In the proof of \cref{cor: X is F-rational}, we identified the generators of $\cO_{\mathfrak{X}}(D_+(W))$ as an $A_a$-algebra to be 
    \[
    \left\{x_1^{a_1}x_2^{a_2}x_3^{a_3}\frac{T_a}{z}\mid a_1+a_2+a_3 = 3\right\} \bigcup \left\{y^3\left(\frac{T_a}{z}\right)^2\right\}\bigcup\left\{yx_i\frac{T_a}{z}\mid 1\leq i\leq 3\right\}.
    \]
    
    If $b\in\Lambda$ then
    \[
    t_b = \frac{T_bz^2}{T_az^2}\quad \mbox{and}\quad T_bz^2\in V_1.
    \]
    If $a_1+a_2+a_3 +2b = 3$, then 
    \[
    x_1^{a_1}x_2^{a_2}x_3^{a_3}y^b\frac{T_a}{z} = \frac{T_a^2zx_1^{a_1}x_2^{a_2}x_3^{a_3}y^b}{T_az^2} \quad \mbox{and}\quad T_a^2zx_1^{a_1}x_2^{a_2}x_3^{a_3}y^b\in V_2.
    \]
    The final generator is
    \[
    y^3\left(\frac{T_a}{z}\right)^2=\frac{T_a^3y^3}{T_az^2}\quad\mbox{and}\quad T_a^3y^3\in V_3,
    \]
    so
    \[
    \cO_\mathfrak{X}(D_+(T_az^2)) = k\left[\frac{U}{T_az^2}\mid U\in V\right]
    \]
    as needed.

    We now show that if $a\in \Lambda$, then $\cO_{\mathfrak{X}}(D_+(T_a^3y^3))$ is a $k$-algebra generated by the sections $\left\{\frac{U}{T_a^3y^3}\mid U\in V\right\}$. We again use the $A_a$-algebra generators of $\cO_{\mathfrak{X}}(D_+(T_a^3y^3))$ from the proof of \cref{cor: X is F-rational}, 
    \[
    \left\{\frac{x_ix_j}{y}\mid 1\leq i\leq j\leq3\right\}\bigcup \left\{\frac{zx_i}{T_ay^2}\mid 1\leq i\leq 3\right\} \bigcup \left\{\frac{z^2}{T_a^2y^3}\right\}.
    \]
    If $b\in\Lambda$, then
    \[
    t_b = \frac{T_a^2T_by^3}{T_a^3y^3}\quad\mbox{and}\quad T_a^2T_by^3\in V_3.
    \]
    If $1\leq i\leq j\leq 3$ then
    \[
    \frac{x_ix_j}{y} = \frac{T_a^3x_ix_jy^2}{T_a^3y^3}\quad\mbox{and}\quad T_a^3x_ix_jy^2\in V_3.
    \]
    If $1\leq i\leq 3$ then
    \[
    \frac{zx_i}{T_ay^2} = \frac{T_a^2zyx_i}{T_a^3y^3}\quad \mbox{and}\quad T_a^2zyx_i\in V_2.
    \]
    The final generator
    \[
    \frac{z^2}{T_a^2y^3} = \frac{T_az^2}{T_a^3y^3}\quad\mbox{and}\quad T_az^2\in V_1.
    \]
    This completes the argument that
    \[
    \cO_\mathfrak{X}(D_+(T_a^3y^3)) = k\left[\frac{U}{T_a^3y^3}\mid U\in V\right]
    \]

    The final class of charts we need to check is a chart of the form 
    \[
    \cO_\mathfrak{X}(D_+(T_a^3x_i^6)) \cong \left(C_a\left[z, T_a,T_a^{-1},x_i^{-6}\right]\right)_{(0,0)} = \left(C_a\left[z, T_a,T_a^{-1},x_i^{-1}\right]\right)_{(0,0)}.
    \]
    Again by the proof of \cref{cor: X is F-rational}, we have $A_a$-algebra generators of $\cO_\mathfrak{X}(D_+(T_a^3x_i^6))$ given by
    \[
    \left\{\frac{x_j}{x_i}\mid j\not=i\right\}\bigcup\left\{\frac{y}{x_i^2}, \frac{z}{T_ax_i^3}\right\}.
    \]
    If $b\in\Lambda$ then
    \[
    t_b = \frac{T_a^2T_bx_i^6}{T_a^3x_i^6}\quad\mbox{and}\quad T_a^2T_bx_i^6\in V_3.
    \]
    If $j\not=i$,
    \[
    \frac{x_j}{x_i} = \frac{T_a^3x_i^5x_j}{T_a^3x_i^6}\quad\mbox{and}\quad T_a^3x_i^5x_j\in V_3.
    \]
    As for the final two generators
    \[
    \frac{y}{x_i^2} = \frac{T_a^3x_i^4y}{T_a^3x_i^6}\quad\mbox{and}\quad T_a^3x_i^4y\in V_3
    \]
    and
    \[
    \frac{z}{T_ax_i^3} = \frac{T_a^2zx_i^3}{T_a^3x_i^6}\quad\mbox{and}\quad T_a^2zx_i^3\in V_2.
    \]
    Therefore $\cO_\mathfrak{X}(D_+(T_a^3x_i^6)) = k\left[\frac{U}{T_a^3x_i^6}\mid U\in V\right]$ as needed to complete the proof.  
\end{proof}

\end{document}
